% problem-000406 5 化学反应与结构综合分析
\section*{5. 化学反应与结构综合分析}
（图：）
（THF为四氢呋喃，HMPA为六甲基磷酸三胺，是常⽤有机溶剂；+BuOH为叔丁醇）

解答以下问题时不考虑反应过程中的⽴体化学。

\subsection*{5-1}
写出A、B、C的结构和D代表的反应条件。

\subsection*{5-2}
化合(5)在叔丁醇中加热处理，依次经3个中间物E、F、G，⽣成产物(6)，E和G各含3个羰基，F含2个羰基。

\subsection*{5-2}
-1 请写出E、F、G的结构。

\subsection*{5-2}
-2 ⽤电⼦转移箭头标明G到(6)的机理。

\subsection*{5-3}
⽤格⽒试剂(7)和化合物(1)反应，然后依次⽤CH_{3}I、 $\mathrm { C H _ { 3 } C O O H }$ 处理，得到化合物(8)，如下式所⽰：

（图：）

请写出从J到(8)的反应历程，⽤弯箭头标明电⼦转移的⽅向，并说明(1)与(2)和(1)与(7)反应，为何采⽤相同的反应条件，却得到不同类别的反应产物(3)和(8)（⽤结构简式表⽰）。
